\problemname{Rarest Insects}

\begin{quote}
\textit{Note:} in this problem, you have to write your solution in C++.
\end{quote}

There are $N$ insects, indexed from $0$ to $N-1$, running around Pak Blangkon's house. Each insect has
a type, which is an integer between $0$ and $10^9$ inclusive. Multiple insects may have the same type.

Suppose insects are grouped by type. We define the \emph{cardinality of the most frequent insect type}
as the number of insects in a group with the most number of insects. Similarly, the \emph{cardinality
of the rarest insect type} is the number of insects in a group with the least number of insects.

For example, suppose that there are $11$ insects, whose types are $[5, 7, 9, 11, 11, 5, 0, 11, 9, 100,
9]$. In this case, the cardinality of the most frequent insect type is $3$. The groups with the most
number of insects are type $9$ and type $11$, each consisting of $3$ insects. The cardinality of the
rarest insect type is $1$. The groups with the least number of insects are type $7$, type $0$, and type
$100$, each consisting of $1$ insect.

Pak Blangkon does not know the type of any insect. He has a machine with a single button that can
provide some information about the types of the insects. Initially, the machine is empty. To use the
machine, three types of operations can be performed:
\begin{enumerate}
  \item Move an insect to inside the machine.
  \item Move an insect to outside the machine.
  \item Press the button on the machine.
\end{enumerate}

Each type of operation can be performed at most $40\,000$ times.

Whenever the button is pressed, the machine reports the cardinality of the most frequent insect type,
considering only insects inside the machine.

Your task is to determine the cardinality of the rarest insect type among all $N$ insects in Pak
Blangkon's house by using the machine. Additionally, in test group $3$ your score depends on the
maximum number of calls of a given type that are performed; see the Scoring section for details.

\section*{Implementation Details}
You should implement the following procedure:

\begin{verbatim}
int min_cardinality(int N)
\end{verbatim}

\begin{itemize}
  \item $N$: the number of insects.
  \item This procedure should return the cardinality of the rarest insect type among all $N$ insects
  in Pak Blangkon's house.
  \item This procedure is called exactly once.
\end{itemize}

The above procedure can make calls to the following procedures:

\begin{verbatim}
void move_inside(int i)
\end{verbatim}

\begin{itemize}
  \item $i$: the index of the insect to be moved inside the machine. The value of $i$ must be between
  $0$ and $N-1$ inclusive.
  \item If this insect is already inside the machine, the call has no effect on the set of insects in
  the machine. However, it is still counted as a separate call.
  \item This procedure can be called at most $40\,000$ times.
\end{itemize}

\begin{verbatim}
void move_outside(int i)
\end{verbatim}

\begin{itemize}
  \item $i$: the index of the insect to be moved outside the machine. The value of $i$ must be
  between $0$ and $N-1$ inclusive.
  \item If this insect is already outside the machine, the call has no effect on the set of insects in
  the machine. However, it is still counted as a separate call.
  \item This procedure can be called at most $40\,000$ times.
\end{itemize}

\begin{verbatim}
int press_button()
\end{verbatim}

\begin{itemize}
  \item This procedure returns the cardinality of the most frequent insect type, considering only
  insects inside the machine.
  \item This procedure can be called at most $40\,000$ times.
\end{itemize}

If you call one of the three procedures more often than that, pass an $i$ that is not between $0$ and
$N-1$, or return an incorrect value from \texttt{min\_cardinality}, you will receive $0$ points for
the test group.

Your submission must not contain a \texttt{main} function, and must not read from standard input or
write to standard output: a grader is compiled together with your submission and takes care of all
communication with the machine. The declarations above live in \texttt{insects.h}, which is available
both to your submission and as an attachment, so start your file with
\texttt{\#include "insects.h"}.

The judge is \textbf{not adaptive}. That is, the types of all $N$ insects are fixed before
\texttt{min\_cardinality} is called.

\section*{Constraints}
\begin{itemize}
  \item $2 \leq N \leq 2000$
\end{itemize}

\section*{Scoring}
Your solution will be tested on a set of test groups.
To earn points for a group, you must pass all test cases in that group.

\noindent
\begin{tabular}{| l | l | p{12cm} |}
  \hline
  \textbf{Group} & \textbf{Points} & \textbf{Constraints} \\ \hline
  $1$ & $10$ & $N \leq 200$ \\ \hline
  $2$ & $15$ & $N \leq 1000$ \\ \hline
  $3$ & $75$ & No additional constraints. \\ \hline
\end{tabular}

\medskip
\noindent
Let $q$ be the maximum of the following three values: the number of calls to
\texttt{move\_inside}, the number of calls to \texttt{move\_outside}, and the number of calls to
\texttt{press\_button}.

In group $3$ you can obtain a partial score. Let $m$ be the maximum value of $\frac{q}{N}$ across all
test cases in that group. Your score for group $3$ is calculated according to the following table:

\begin{center}
\begin{tabular}{| l | l |}
  \hline
  \textbf{Condition} & \textbf{Points} \\ \hline
  $20 < m$          & $0$ \\ \hline
  $6 < m \leq 20$   & $\frac{225}{m-2}$ \\ \hline
  $3 < m \leq 6$    & $81 - \frac{2}{3} m^2$ \\ \hline
  $m \leq 3$        & $75$ \\ \hline
\end{tabular}
\end{center}

\noindent
Points are rounded to two decimals.

\section*{Sample Grader}
A sample grader is available under ``attachments'' at the bottom of this page, together with
\texttt{insects.h} and the sample input. Compile it with your solution and run it on an input file to
try your solution locally. Let $T$ be an array of $N$ integers where $T[i]$ is the type of insect
$i$. The sample grader reads the input in the following format:
\begin{itemize}
  \item line $1$: $N$
  \item line $2$: $T[0]$ $T[1]$ $\ldots$ $T[N-1]$
\end{itemize}

If the sample grader detects that your solution breaks the rules, its output is
\texttt{Protocol Violation: <MSG>}, where \texttt{<MSG>} is one of the following:
\begin{itemize}
  \item \texttt{invalid parameter}: in a call to \texttt{move\_inside} or \texttt{move\_outside}, the
  value of $i$ is not between $0$ and $N-1$ inclusive.
  \item \texttt{too many calls}: the number of calls to any of \texttt{move\_inside},
  \texttt{move\_outside} or \texttt{press\_button} exceeds $40\,000$.
\end{itemize}

Otherwise, the output of the sample grader is in the following format:
\begin{itemize}
  \item line $1$: the return value of \texttt{min\_cardinality}
  \item line $2$: $q$
\end{itemize}

The sample grader runs the machine inside your own process; the judge runs it as a separate program.

That is the format of the sample data below. The sample output is what the sample grader prints for
the sequence of calls in the table below; since $q$ counts \emph{your} calls, the second line will
differ for a solution that queries the machine differently.

\section*{Explanation of Sample}
The sample is a scenario in which there are $6$ insects of types $[5, 8, 9, 5, 9, 9]$ respectively.
The procedure \texttt{min\_cardinality} is called with $N = 6$, and may call \texttt{move\_inside},
\texttt{move\_outside} and \texttt{press\_button} as follows:

\begin{center}
\begin{tabular}{| l | l | l | l |}
  \hline
  \textbf{Call} & \textbf{Return value} & \textbf{Insects in the machine} &
  \textbf{Types of insects in the machine} \\ \hline
  & & $\{\}$ & $[\,]$ \\ \hline
  \texttt{move\_inside(0)}  &     & $\{0\}$           & $[5]$ \\ \hline
  \texttt{press\_button()}  & $1$ & $\{0\}$           & $[5]$ \\ \hline
  \texttt{move\_inside(1)}  &     & $\{0, 1\}$        & $[5, 8]$ \\ \hline
  \texttt{press\_button()}  & $1$ & $\{0, 1\}$        & $[5, 8]$ \\ \hline
  \texttt{move\_inside(3)}  &     & $\{0, 1, 3\}$     & $[5, 8, 5]$ \\ \hline
  \texttt{press\_button()}  & $2$ & $\{0, 1, 3\}$     & $[5, 8, 5]$ \\ \hline
  \texttt{move\_inside(2)}  &     & $\{0, 1, 2, 3\}$  & $[5, 8, 9, 5]$ \\ \hline
  \texttt{move\_inside(4)}  &     & $\{0, 1, 2, 3, 4\}$ & $[5, 8, 9, 5, 9]$ \\ \hline
  \texttt{move\_inside(5)}  &     & $\{0, 1, 2, 3, 4, 5\}$ & $[5, 8, 9, 5, 9, 9]$ \\ \hline
  \texttt{press\_button()}  & $3$ & $\{0, 1, 2, 3, 4, 5\}$ & $[5, 8, 9, 5, 9, 9]$ \\ \hline
  \texttt{move\_inside(5)}  &     & $\{0, 1, 2, 3, 4, 5\}$ & $[5, 8, 9, 5, 9, 9]$ \\ \hline
  \texttt{press\_button()}  & $3$ & $\{0, 1, 2, 3, 4, 5\}$ & $[5, 8, 9, 5, 9, 9]$ \\ \hline
  \texttt{move\_outside(5)} &     & $\{0, 1, 2, 3, 4\}$ & $[5, 8, 9, 5, 9]$ \\ \hline
  \texttt{press\_button()}  & $2$ & $\{0, 1, 2, 3, 4\}$ & $[5, 8, 9, 5, 9]$ \\ \hline
\end{tabular}
\end{center}

\noindent
At this point, there is sufficient information to conclude that the cardinality of the rarest insect
type is $1$. Therefore, \texttt{min\_cardinality} should return $1$.

In this example, \texttt{move\_inside} is called $7$ times, \texttt{move\_outside} is called $1$ time
and \texttt{press\_button} is called $6$ times.
